\documentclass[11pt]{amsart}

\usepackage[T1]{fontenc}
\usepackage[utf8]{inputenc}

\usepackage{amsmath,amssymb,amsthm,mathtools}
\usepackage{enumitem}

\usepackage[colorlinks=true,linkcolor=blue,citecolor=blue,urlcolor=blue]{hyperref}

\setlist[enumerate]{leftmargin=2.2em}
\setlist[itemize]{leftmargin=2.0em}
\allowdisplaybreaks

\newtheorem{theorem}{Theorem}[section]
\newtheorem{proposition}[theorem]{Proposition}
\newtheorem{lemma}[theorem]{Lemma}
\newtheorem{corollary}[theorem]{Corollary}
\newtheorem{claim}[theorem]{Claim}
\theoremstyle{definition}
\newtheorem{definition}[theorem]{Definition}
\newtheorem{remark}[theorem]{Remark}
\newtheorem{example}[theorem]{Example}

\newcommand{\F}{\mathbb F}
\newcommand{\N}{\mathbb N}
\newcommand{\Z}{\mathbb Z}
\newcommand{\Sym}{\operatorname{Sym}}
\newcommand{\GL}{\operatorname{GL}}
\newcommand{\EL}{\operatorname{EL}}
\newcommand{\St}{\operatorname{St}}
\newcommand{\Fix}{\operatorname{Fix}}
\newcommand{\id}{\mathrm{id}}
\newcommand{\1}{\mathbf 1}
\newcommand{\cP}{\mathcal P}
\newcommand{\cB}{\mathcal B}
\newcommand{\cG}{\mathcal G}
\newcommand{\cE}{\mathcal E}
\newcommand{\eps}{\varepsilon}
\newcommand{\dd}{\,\mathrm d}
\newcommand{\distH}{d_{\mathrm H}}
\newcommand{\into}{\hookrightarrow}
\newcommand{\onto}{\twoheadrightarrow}
\newcommand{\abs}[1]{\lvert#1\rvert}
\newcommand{\norm}[1]{\lVert#1\rVert}
\newcommand{\set}[1]{\left\{#1\right\}}
\newcommand{\angles}[1]{\left\langle#1\right\rangle}
\newcommand{\suchthat}{\,\middle|\,}
\newcommand{\symdiff}{\mathbin{\triangle}}
\newcommand{\oN}{o(N_n)}
\newcommand{\on}{o_n(1)}

\DeclareMathOperator{\med}{med}
\DeclareMathOperator{\supp}{supp}
\DeclareMathOperator{\diag}{diag}

\title{Nonsofic Groups Exist}
\author{Anonymous}
% Replace the author block before submission.
\date{August 2026}

\subjclass[2020]{20F65, 20E26, 20F05, 05C48, 16S88}
\keywords{sofic groups, property (T), expander graphs, Leavitt algebras, Thompson's group}

\begin{document}

\begin{abstract}
We prove that the group \(\EL_4\bigl(L_{\F_2}(1,2)\bigr)\) of elementary matrices over the binary Leavitt algebra is not sofic.  The key analytic input is an expander--matching criterion: if a property-\((T)\) group is generated by an infinite property-\((T)\) subgroup \(\Gamma\) together with finitely many elements that conjugate \(\Gamma\) into itself, then in a hypothetical sofic approximation the component decompositions supplied by Kun's theorem for \(\Gamma\) and for the ambient group can be matched almost bijectively.  The matching combines a median normalization of component sizes, exact conservation under permutations, and an \(\ell^1\) co-area estimate on the ambient expanders.  Restriction to one matched component produces a sofic approximation of \(K\times J\), where \(K\) is the compressed image of \(\Gamma\) and \(J\) is a commuting subgroup, in which the \(K\)-edges form an essential expander sequence; the Kun--Thom centralizer obstruction then forces \(J\) to be LEF.  For the displayed group we exhibit two explicit compressors with a common image and a commuting copy of Thompson's group \(V\), which is infinite, simple, and finitely presented, hence not LEF.  Every invertible matrix used is written as an explicit product of elementary matrices, so no \(K\)-theory enters.  Finally, a finite multiplication-table presentation yields an infinite finitely presented nonsofic group, which can moreover be chosen to have property~\((T)\).
\end{abstract}

\maketitle
\tableofcontents

\section{Introduction}

A countable discrete group is \emph{sofic} if its finite multiplication tables admit arbitrarily accurate and asymptotically faithful realizations by permutations of finite sets.  The approximation property arose in Gromov's work on symbolic dynamics \cite{Gromov99}; Weiss introduced the term \emph{sofic} and asked whether every countable group is sofic \cite{Weiss00}.  The class contains all amenable and all residually finite groups and is stable under a number of standard constructions; see, for example, \cite{ElekSzabo06,Pestov08}.  The question whether every group is sofic has remained open.

We give a negative answer using the binary Leavitt algebra
\begin{equation}\label{eq:Leavitt-presentation}
 R=L_{\F_2}(1,2)
 =\F_2\langle s_0,s_1,t_0,t_1
 \mid t_i s_j=\delta_{ij},\ s_0t_0+s_1t_1=1\rangle.
\end{equation}
For a unital ring \(A\), let \(\EL_m(A)\) denote the subgroup of \(\GL_m(A)\) generated by the elementary matrices
\[
 x_{ij}(a)=I_m+aE_{ij},\qquad a\in A,\quad i\ne j.
\]
Our first example is the Kazhdan group \(\EL_4(R)\).  An infinite finitely presented example is then extracted from a finite part of its multiplication table.

\begin{theorem}[Main theorem]\label{thm:main}
Let \(R=L_{\F_2}(1,2)\).  The group
\[
 G=\EL_4(R)
\]
is infinite, finitely generated, has property~\((T)\), and is not sofic.  Moreover, for a suitable finite subset \(F\subseteq G\), the group \(H_F\) defined in \eqref{eq:HF-presentation} below is infinite, finitely presented, and not sofic, and it may be replaced by an infinite finitely presented group with property~\((T)\) that is not sofic.  In particular, not every finitely presented discrete group is sofic.
\end{theorem}

The main analytic issue is the distinction between a single expander and a disjoint union of expanders.  Kun proved that a sofic approximation of a property-\((T)\) group can be changed on a vanishing proportion of edges into a disjoint union of uniformly expanding components \cite{Kun19}.  Kun and Thom proved that if the restriction of a sofic approximation of \(\Lambda\times\Delta\) to a Kazhdan factor \(\Lambda\) is essentially one expander, then \(\Delta\) is locally embeddable into finite groups (LEF) \cite{KunThom19}.  A disjoint union of equal-sized expanders is not enough: the commuting group may permute the components.

The Leavitt relations provide a way to pass from the disjoint union to one component.  We construct a property-\((T)\) subgroup
\[
 \Gamma\cong\EL_3(R)
\]
and two elements \(u,v\in G\) such that
\[
 G=\langle\Gamma,u,v\rangle,
 \qquad
 u\Gamma u^{-1}=v\Gamma v^{-1}=:K\leq\Gamma.
\]
A hypothetical sofic approximation gives a Kun decomposition for \(\Gamma\) and another for \(G\).  Transporting the \(\Gamma\)-components by a compressor gives one-sided inclusions into other \(\Gamma\)-components.  On each ambient \(G\)-component, normalize a \(\Gamma\)-component size \(s\) by
\[
 f=\frac{s}{s+m},
\]
where \(m\) is a vertex-weighted median of the component sizes.  The one-sided inclusions imply almost monotonicity of \(f\) along compressor arcs.  Since each compressor is a permutation, total increase equals total decrease.  The ambient expander co-area estimate then forces \(f\) to concentrate at its median value \(1/2\).  Hence the transported and target components have asymptotically equal sizes, upgrading one-sided inclusion to almost-bijection.

We restrict to one well-chosen matched component.  The transported copy of \(\Gamma\) gives a single essentially expanding approximation of \(K\), while every element of a commuting subgroup \(J\leq\Gamma\) almost preserves the target component.  Completing the restrictions to permutations produces a sofic approximation of \(K\times J\) to which Kun--Thom applies.

Algebraically, \(J\) is a copy of Thompson's group \(V\), placed in the complementary binary corner.  The group \(V\) is infinite, simple, and finitely presented \cite{Higman74,BleakQuick17}; it is therefore not LEF.  The resulting contradiction proves nonsoficity of \(G\).

The final finite-presentation argument does not claim that soficity passes to quotients.  Nonsoficity supplies a finite multiplication-table test that no finite permutation model can satisfy.  Imposing exactly that finite table gives an infinite finitely presented group that retains the obstruction.  Shalom's finitely presented covers of Kazhdan groups \cite{Shalom00} then supply a version with property~\((T)\).

\subsection*{External inputs}

The proof imports two results on sofic approximations of Kazhdan groups: Kun's expander decomposition theorem (Theorem~\ref{thm:Kun}, from \cite{Kun19}) and the Kun--Thom centralizer obstruction (Theorem~\ref{thm:Kun-Thom}, from \cite{KunThom19}); both are stated in Section~\ref{sec:preliminaries} in the form in which they are used.  The only further external inputs are property~\((T)\) for elementary groups over finitely generated rings \cite{ErshovJaikin10}, Shalom's finitely presented Kazhdan covers \cite{Shalom00}, and classical facts about Thompson's group \(V\) \cite{Higman74,CannonFloydParry96,BleakQuick17}.  Everything else is self-contained.  In particular, every invertible matrix used is written as an explicit product of elementary matrices, so no \(K\)-theory enters, and the analytic argument rests on an elementary \(\ell^1\) co-area estimate rather than on spectral theory.

\section{Sofic approximations, expanders, and LEF groups}\label{sec:preliminaries}

\subsection{Sofic approximations}

For a nonempty finite set \(Y\), the normalized Hamming distance on \(\Sym(Y)\) is
\[
 \distH(\sigma,\tau)
 =\frac{1}{\abs{Y}}\abs{\set{y\in Y\suchthat \sigma(y)\ne\tau(y)}}.
\]

\begin{definition}\label{def:sofic-approx}
A sequence of maps
\[
 \sigma_n:G\longrightarrow\Sym(Y_n)
\]
is a \emph{sofic approximation} of a countable group \(G\) if, for all \(g,h\in G\),
\[
 \distH\bigl(\sigma_n(g)\sigma_n(h),\sigma_n(gh)\bigr)\longrightarrow0,
\]
and for every \(g\ne1\),
\[
 \distH\bigl(\sigma_n(g),\id_{Y_n}\bigr)\longrightarrow1.
\]
The group \(G\) is \emph{sofic} if it admits a sofic approximation.
\end{definition}

For every fixed finite subset of \(G\), changing each \(\sigma_n(g)\) on \(o(\abs{Y_n})\) points does not affect the approximation.  We shall therefore freely make the following standard normalization.

\begin{lemma}[Finite normalization]\label{lem:finite-normalization}
Let \(E=E^{-1}\subseteq G\) be finite and contain \(1\).  After changing \(\sigma_n(g)\) on \(o(\abs{Y_n})\) points for \(g\in E\), one may assume
\[
 \sigma_n(1)=\id,
 \qquad
 \sigma_n(g^{-1})=\sigma_n(g)^{-1}
 \quad(g\in E).
\]
\end{lemma}

\begin{proof}
The sofic identities give
\[
 \distH(\sigma_n(1),\id)\to0
 \quad\text{and}\quad
 \distH\bigl(\sigma_n(g^{-1}),\sigma_n(g)^{-1}\bigr)\to0,
\]
the second by bi-invariance of \(\distH\).  Set \(\sigma_n(1)=\id\).  From each pair \(\{g,g^{-1}\}\subseteq E\) with \(g^2\ne1\), keep the permutation of one element and redefine the permutation of the other to be its inverse.  If \(g=g^{-1}\ne1\), replace \(\sigma_n(g)\) by the involution agreeing with \(\sigma_n(g)\) on the fixed points and genuine \(2\)-cycles of \(\sigma_n(g)\) and equal to the identity elsewhere; the modified points lie in \(\{y:\sigma_n(g)^2y\ne y\}\), a set of size \(o(\abs{Y_n})\).  Only finitely many elements are involved, and each change is on \(o(\abs{Y_n})\) points, so the sofic identities persist.
\end{proof}

Let \(S=S^{-1}\) be a finite generating set.  The \(S\)-generator graph associated with \(\sigma_n\) has vertex set \(Y_n\) and a directed arc
\[
 y\xrightarrow{s}\sigma_n(s)y
\]
for each \(s\in S\).  After the normalization above, forgetting orientation and labels gives a bounded-degree undirected multigraph.  All edge-counting estimates below are insensitive to loops, multiplicities, and bounded factors coming from orientation.

\subsection{Expanders and essential expansion}

For a finite graph \(X\) and \(A\subseteq V(X)\), let \(\partial_X A\) be the set of edges with one endpoint in \(A\) and one outside.  The Cheeger constant is
\[
 h(X)=\min_{0<\abs{A}\leq\abs{V(X)}/2}
 \frac{\abs{\partial_X A}}{\abs{A}}.
\]
A bounded-degree sequence \((X_n)\) is an \emph{expander sequence} if \(\abs{V(X_n)}\to\infty\) and \(\inf_n h(X_n)>0\).

We say that two bounded-degree graph sequences on the same vertex sets are \emph{essentially equal} if the symmetric difference of their edge multisets has size \(o(\abs{V(X_n)})\).  A sequence is an \emph{essential expander sequence} if it is essentially equal to a bounded-degree expander sequence.  Edge multisets are used throughout, so loops and multiplicities cause no ambiguity.

Kun's theorem \cite{Kun19} proves Bowen's conjecture: every sofic approximation of a countably infinite property-\((T)\) group is essentially a vertex-disjoint union of expander graphs.  We use it in the following form; the bounded-degree conclusion is part of Kun's formulation.

\begin{theorem}[Kun \cite{Kun19}]\label{thm:Kun}
Let \(\Lambda\) be an infinite finitely generated property-\((T)\) group and let \(S\) be a finite symmetric generating set.  For every sofic approximation of \(\Lambda\), the associated \(S\)-generator graphs are essentially equal to graphs of uniformly bounded degree whose connected components have Cheeger constants bounded below by one positive constant independent of the component and of \(n\).
\end{theorem}

Uniform Cheeger expansion gives a uniform \(\ell^1\) Poincar\'e inequality around any median.  We record the form needed here; it is a discrete co-area argument and requires no spectral theory.

\begin{lemma}[\(\ell^1\) co-area estimate]\label{lem:poincare}
Let \(X\) be a finite graph with \(h(X)\geq h>0\), let \(f:V(X)\to\mathbb R\), and let \(c\in\mathbb R\) satisfy
\[
 \abs{\{f>c\}}\leq\tfrac12\abs{V(X)}
 \qquad\text{and}\qquad
 \abs{\{f<c\}}\leq\tfrac12\abs{V(X)}.
\]
Then
\[
 \sum_{x\in V(X)}\abs{f(x)-c}
 \leq\frac1h\sum_{\{x,y\}\in E(X)}\abs{f(x)-f(y)}.
\]
\end{lemma}

\begin{proof}
For \(t\in\mathbb R\) let \(F_t=\{x:f(x)>t\}\).  An edge \(\{x,y\}\) with \(f(y)<f(x)\) lies in \(\partial_XF_t\) exactly for \(t\in[f(y),f(x))\), so
\[
 \sum_{\{x,y\}\in E(X)}\abs{f(x)-f(y)}
 =\int_{-\infty}^{\infty}\abs{\partial_XF_t}\dd t .
\]
For \(t\geq c\) we have \(\abs{F_t}\leq\abs{V(X)}/2\), hence \(\abs{\partial_XF_t}\geq h\abs{F_t}\), and
\(\int_c^\infty\abs{F_t}\dd t=\sum_x(f(x)-c)_+\).
For \(t<c\), the complement satisfies \(\abs{V(X)\setminus F_t}=\abs{\{f\leq t\}}\leq\abs{\{f<c\}}\leq\abs{V(X)}/2\), hence \(\abs{\partial_XF_t}\geq h\abs{V(X)\setminus F_t}\), and
\(\int_{-\infty}^{c}\abs{\{f\leq t\}}\dd t=\sum_x(c-f(x))_+\).
Adding the two estimates proves the lemma.
\end{proof}

\subsection{LEF groups and the Kun--Thom obstruction}

\begin{definition}
A group \(D\) is \emph{locally embeddable into finite groups}, abbreviated \emph{LEF} \cite{VershikGordon97}, if for every finite \(F\subseteq D\) there are a finite group \(H\) and an injective map \(\phi:F\to H\) such that
\[
 \phi(x)\phi(y)=\phi(xy)
\]
whenever \(x,y,xy\in F\).
\end{definition}

\begin{lemma}\label{lem:fp-lef-rf}
Every finitely presented LEF group is residually finite.
\end{lemma}

\begin{proof}
Write \(D=\langle S\mid r_1,\dots,r_k\rangle\), and let a word \(w\) in \(S^{\pm1}\) represent \(1\ne g\in D\).  Let \(F\subseteq D\) contain \(1\), the elements of \(S^{\pm1}\), and the elements represented by every prefix of each \(r_i\) and of \(w\); enlarge \(F\), still finitely, so that each multiplication used to pass from one such prefix to the next has all three terms in \(F\).

Let \(\phi:F\to H\) be a local embedding into a finite group.  The relations involving \(1\) and inverse letters give \(\phi(1)=1_H\) and \(\phi(s^{-1})=\phi(s)^{-1}\).  Reading each relator one letter at a time shows that the assignment \(s\mapsto\phi(s)\) satisfies every defining relation, and hence induces a homomorphism \(D\to H\).  The same prefix calculation sends \(g\) to \(\phi(g)\ne\phi(1)\), by injectivity of \(\phi\) on \(F\).  Thus every nonidentity element survives in a finite quotient.
\end{proof}

The following is the centralizer obstruction of Kun and Thom \cite{KunThom19}.  Although the statement of their theorem is phrased for an expander sequence, the proof is carried out for an essentially expanding sequence; we record that form.

\begin{theorem}[Kun--Thom \cite{KunThom19}]\label{thm:Kun-Thom}
Let \(\Lambda\) be a countable property-\((T)\) group and let \(D\) be finitely generated.  Suppose that \(\Lambda\times D\) has a sofic approximation for which the edges labelled by a fixed finite generating set of \(\Lambda\) form an essential expander sequence.  Then \(D\) is LEF.
\end{theorem}

\section{An expander--matching criterion}\label{sec:criterion}

We now prove the analytic criterion used in the construction.  The shape of the argument: Kun's theorem applied to \(\Gamma\) supplies expanding source components, and applied to \(G\) it supplies the ambient components on which sizes equilibrate.  Compression gives a one-sided refinement of transported components; permutation conservation makes the comparison two-sided; the co-area estimate pins the normalized sizes at \(1/2\); a weighted selection extracts one matched component; and the Kun--Thom obstruction converts the resulting single expander into the LEF conclusion.

\begin{theorem}[Expander--matching criterion]\label{thm:criterion}
Let \(G\) be a finitely generated property-\((T)\) group, and let \(\Gamma\leq G\) be an infinite finitely generated property-\((T)\) subgroup.  Suppose there is a finite set \(Q\subseteq G\) such that
\begin{equation}\label{eq:criterion-generation}
 G=\langle\Gamma,Q\rangle
\end{equation}
and
\begin{equation}\label{eq:criterion-compression}
 q\Gamma q^{-1}\leq\Gamma
 \qquad(q\in Q).
\end{equation}
Fix \(q_0\in Q\) and put
\[
 K=q_0\Gamma q_0^{-1}.
\]
Let \(J\leq\Gamma\) be finitely generated and suppose
\begin{equation}\label{eq:criterion-centralizer}
 [K,J]=1,
 \qquad
 K\cap J=\{1\}.
\end{equation}
If \(G\) is sofic, then \(J\) is LEF.
\end{theorem}

Throughout the proof, write \(N_n=\abs{Y_n}\), and use \(o(N_n)\) for quantities whose ratio to \(N_n\) tends to zero.  All implicit constants depend only on fixed generating sets.

\subsection{The two component decompositions}

Assume that
\[
 \sigma_n:G\longrightarrow\Sym(Y_n)
\]
is a sofic approximation.  Choose a finite symmetric generating set \(S=S^{-1}\) of \(\Gamma\).  By \eqref{eq:criterion-generation},
\[
 A=S\cup Q\cup Q^{-1}
\]
is a finite symmetric generating set of \(G\); the set \(Q\) itself is not enlarged, since the compression hypothesis \eqref{eq:criterion-compression} is assumed only for \(q\in Q\).  Put \(S_K=q_0Sq_0^{-1}\).  Apply Lemma~\ref{lem:finite-normalization} to the finite symmetric set \(A\cup S_K\cup\{1\}\).

Apply Theorem~\ref{thm:Kun} to the restriction to \(\Gamma\).  We obtain a graph \(\widetilde X_n\) on \(Y_n\), differing from the \(S\)-generator graph in \(o(N_n)\) edges, of maximum degree at most a fixed \(d_\Gamma\), and whose connected components have Cheeger constant at least \(h_\Gamma>0\).  Let
\[
 \cP_n
\]
be its component partition, and let \(P_n(y)\) denote the block containing \(y\).

Apply Theorem~\ref{thm:Kun} to \(G\).  We obtain a graph \(\widetilde Z_n\), differing from the \(A\)-generator graph in \(o(N_n)\) edges, of uniformly bounded degree, and whose components have Cheeger constant at least \(h_G>0\).  Let
\[
 \cB_n
\]
be its component partition, and write \(B_n(y)\) for the block containing \(y\).

For every fixed \(\gamma\in\Gamma\),
\begin{equation}\label{eq:Gamma-block-invariance}
 \abs{\set{y\in Y_n\suchthat
 P_n(\sigma_n(\gamma)y)\ne P_n(y)}}=o(N_n).
\end{equation}
Indeed, for \(\gamma\in S\), every crossing arc must be among the \(o(N_n)\) modified generator edges.  If \(\gamma=s_1\cdots s_\ell\) is a fixed \(S\)-word, then outside the union of those crossing sets and the finitely many multiplication-defect sets, the successive points
\(y,\sigma_n(s_\ell)y,\dots,\sigma_n(s_1)\cdots\sigma_n(s_\ell)y\)
remain in one block and the last point equals \(\sigma_n(\gamma)y\); this proves the estimate by a union bound.  Similarly, for every \(a\in A\),
\begin{equation}\label{eq:G-block-invariance}
 \abs{\set{y\in Y_n\suchthat
 B_n(\sigma_n(a)y)\ne B_n(y)}}=o(N_n).
\end{equation}

\subsection{A refinement estimate}

\begin{lemma}[Expander refinement]\label{lem:refinement}
Let \(X\) be a finite graph with \(h(X)\geq h>0\), and let \(\mathcal R\) be a partition of \(V(X)\).  Let \(R_X\in\mathcal R\) maximize \(\abs{V(X)\cap R}\), and let \(e_{\mathcal R}(X)\) be the number of edges of \(X\) whose endpoints lie in different \(\mathcal R\)-blocks.  Then
\[
 \abs{V(X)\setminus R_X}
 \leq \frac{2}{h}e_{\mathcal R}(X).
\]
\end{lemma}

\begin{proof}
Put \(D=V(X)\) and \(r=\abs{D\setminus R_X}\).  If \(r\leq\abs D/2\), expansion of \(D\setminus R_X\) gives at least \(hr\) crossing edges.

Suppose \(r>\abs D/2\).  Then every \(\mathcal R\)-piece contained in \(D\) has size at most \(\abs D/2\).  Summing the Cheeger inequality over all such pieces gives at least \(h\abs D\) boundary incidences.  Each crossing edge is counted twice, so
\[
 e_{\mathcal R}(X)\geq\frac{h\abs D}{2}\geq\frac{hr}{2}.
\]
Both cases give the assertion.
\end{proof}

Fix \(q\in Q\).  For each \(C\in\cP_n\), transport the expander \(\widetilde X_n[C]\) by \(\sigma_n(q)\).  This gives an \(h_\Gamma\)-expander on \(\sigma_n(q)C\).  Apart from \(o(N_n)\) total edges, its edges agree with edges labelled by the fixed elements
\[
 qSq^{-1}\subseteq\Gamma.
\]
Indeed, the original Kun modifications contribute \(o(N_n)\) edges, and for each \(s\in S\), sofic multiplicativity gives
\begin{equation}\label{eq:conjugacy-defect}
 \sigma_n(q)\sigma_n(s)\sigma_n(q)^{-1}
 \sim
 \sigma_n(qsq^{-1})
\end{equation}
with Hamming error \(o(1)\).

By \eqref{eq:Gamma-block-invariance}, only \(o(N_n)\) of the transported edges cross the partition \(\cP_n\).  Lemma~\ref{lem:refinement}, summed over the transported components, gives a map
\[
 \Phi_{q,n}:\cP_n\longrightarrow\cP_n
\]
such that, with
\begin{equation}\label{eq:e-q-C}
 e_{q,n}(C)
 =\abs{\sigma_n(q)C\setminus\Phi_{q,n}(C)},
\end{equation}
we have
\begin{equation}\label{eq:sum-e-q-C}
 \sum_{C\in\cP_n}e_{q,n}(C)=o(N_n).
\end{equation}
Thus each transported component is almost contained, in a global weighted sense, in one original \(\Gamma\)-component.

\subsection{Median normalization and conservation}

Define
\[
 s_n(y)=\abs{P_n(y)}.
\]
For every \(B\in\cB_n\), choose a vertex-weighted median \(m_n(B)\) of the values \(s_n(y)\), \(y\in B\): at least half the vertices of \(B\) satisfy \(s_n(y)\leq m_n(B)\), and at least half satisfy \(s_n(y)\geq m_n(B)\).  Define
\begin{equation}\label{eq:median-function}
 f_n(y)=\frac{s_n(y)}{s_n(y)+m_n(B_n(y))}.
\end{equation}
Then \(0<f_n<1\), and \(1/2\) is a median of \(f_n\) on every \(B\in\cB_n\).

For \(s\in S\), if both the \(s\)-arc stays in the same \(\cP_n\)-block and in the same \(\cB_n\)-block, then \(f_n\) is unchanged.  Equations \eqref{eq:Gamma-block-invariance} and \eqref{eq:G-block-invariance} therefore imply
\begin{equation}\label{eq:f-S-variation}
 \sum_{y\in Y_n}
 \abs{f_n(\sigma_n(s)y)-f_n(y)}=o(N_n)
 \qquad(s\in S).
\end{equation}

Now fix \(q\in Q\).  Let \(C=P_n(y)\) and \(D=\Phi_{q,n}(C)\).  Suppose
\begin{equation}\label{eq:good-q-point}
 \sigma_n(q)y\in D
 \quad\text{and}\quad
 B_n(\sigma_n(q)y)=B_n(y)=B.
\end{equation}
Since \(\abs{\sigma_n(q)C}=\abs C\),
\[
 \abs D\geq\abs C-e_{q,n}(C).
\]
For \(a,m>0\) and \(0\leq\delta\leq1\),
\begin{equation}\label{eq:fraction-loss}
 \frac{a}{a+m}
 -\frac{(1-\delta)a}{(1-\delta)a+m}
 \leq\delta.
\end{equation}
Applying this with \(a=\abs C\), \(m=m_n(B)\), and \(\delta=e_{q,n}(C)/\abs C\), we obtain on the points satisfying \eqref{eq:good-q-point}
\[
 \bigl(f_n(y)-f_n(\sigma_n(q)y)\bigr)_+
 \leq\frac{e_{q,n}(C)}{\abs C}.
\]
Summing first over each \(C\), and adding the exceptional points in \eqref{eq:e-q-C} and the \(q\)-arcs crossing \(\cB_n\), gives the explicit bound
\[
 \sum_y\bigl(f_n(y)-f_n(\sigma_n(q)y)\bigr)_+
 \leq 2\sum_{C\in\cP_n}e_{q,n}(C)
 +\abs{\set{y:B_n(\sigma_n(q)y)\ne B_n(y)}}.
\]
Consequently,
\begin{equation}\label{eq:one-sided-variation}
 \sum_y\bigl(f_n(y)-f_n(\sigma_n(q)y)\bigr)_+
 =o(N_n).
\end{equation}

Permutation conservation is exact:
\begin{equation}\label{eq:permutation-conservation}
 \sum_y f_n(\sigma_n(q)y)=\sum_y f_n(y).
\end{equation}
Consequently the total positive and negative variations are equal, and \eqref{eq:one-sided-variation} implies
\begin{equation}\label{eq:f-Q-variation}
 \sum_y\abs{f_n(\sigma_n(q)y)-f_n(y)}=o(N_n)
 \qquad(q\in Q\cup Q^{-1}).
\end{equation}
For \(q^{-1}\), the sum equals the sum for \(q\) after the substitution \(y=\sigma_n(q)x\), using the normalization \(\sigma_n(q^{-1})=\sigma_n(q)^{-1}\).

Equations \eqref{eq:f-S-variation} and \eqref{eq:f-Q-variation} show that the total \(\ell^1\) increment of \(f_n\) over the edges of the \(A\)-generator graph is \(o(N_n)\).  Since \(0<f_n<1\), changing \(o(N_n)\) edges does not affect that conclusion, so the same bound holds over the edge set of \(\widetilde Z_n\).  On every component \(B\in\cB_n\), at most half the vertices satisfy \(f_n>1/2\) and at most half satisfy \(f_n<1/2\), because \(m_n(B)\) is a vertex-weighted median of the sizes.  Lemma~\ref{lem:poincare}, applied with \(c=1/2\) on every component of \(\widetilde Z_n\) and summed over the components, yields
\begin{equation}\label{eq:concentration-half}
 \frac1{N_n}\sum_{y\in Y_n}
 \abs{f_n(y)-\tfrac12}\longrightarrow0.
\end{equation}
The normalized constant \(1/2\) means precisely that component sizes concentrate at the median:
\[
 \frac{s}{s+m}\to\frac12
 \quad\Longleftrightarrow\quad
 \frac{s}{m}\to1.
\]

\subsection{Almost-bijective matching}

\begin{proposition}[Component matching]\label{prop:component-matching}
For every fixed \(q\in Q\), there are collections \(\cG_{q,n}\subseteq\cP_n\) such that
\begin{enumerate}[label=\textup{(\roman*)}]
\item \(\displaystyle\sum_{C\in\cG_{q,n}}\abs C=(1-o_n(1))N_n\);
\item uniformly for \(C\in\cG_{q,n}\),
\begin{equation}\label{eq:symmetric-matching}
 \abs{\sigma_n(q)C\symdiff\Phi_{q,n}(C)}=o_n(1)\abs C;
\end{equation}
\item the restriction of \(\Phi_{q,n}\) to \(\cG_{q,n}\) is injective;
\item the distinct target blocks have total mass \((1-o_n(1))N_n\).
\end{enumerate}
\end{proposition}

\begin{proof}
Choose \(\eta_n\downarrow0\) sufficiently slowly that all global errors in \eqref{eq:G-block-invariance}, \eqref{eq:sum-e-q-C}, and \eqref{eq:concentration-half} are \(o(\eta_n^2N_n)\).  Let \(W_n\) be the union of the following bad vertex sets:
\begin{itemize}
\item points \(y\) with \(\abs{f_n(y)-1/2}>\eta_n\);
\item points \(y\) with \(\abs{f_n(\sigma_n(q)y)-1/2}>\eta_n\);
\item points whose \(q\)-arc crosses \(\cB_n\);
\item points \(y\in C\) for which \(\sigma_n(q)y\notin\Phi_{q,n}(C)\).
\end{itemize}
Equations already proved give \(\abs{W_n}=o(\eta_nN_n)\); for the first two items, combine Markov's inequality with \eqref{eq:concentration-half}.  Discard every source block \(C\) for which
\[
 e_{q,n}(C)>\eta_n\abs C
 \quad\text{or}\quad
 \abs{C\cap W_n}>\eta_n\abs C.
\]
The discarded blocks have total mass \(o(N_n)\).  Call the remaining collection \(\cG_{q,n}\).

For \(C\in\cG_{q,n}\), choose \(y\in C\setminus W_n\), put \(D=\Phi_{q,n}(C)\), and let \(B=B_n(y)=B_n(\sigma_n(q)y)\).  From \eqref{eq:median-function},
\[
 \abs{f_n(y)-\tfrac12}\leq\eta_n
 \quad\Longrightarrow\quad
 \frac{1-2\eta_n}{1+2\eta_n}
 \leq\frac{\abs C}{m_n(B)}
 \leq\frac{1+2\eta_n}{1-2\eta_n}.
\]
The same estimate for \(\sigma_n(q)y\in D\) gives the identical bound with \(\abs D\) in place of \(\abs C\).  Hence
\begin{equation}\label{eq:matched-size-ratio}
 \abs D=(1+O(\eta_n))\abs C.
\end{equation}
Together with \(e_{q,n}(C)\leq\eta_n\abs C\), this gives
\[
 \begin{aligned}
 \abs{\sigma_n(q)C\symdiff D}
 &=\abs{\sigma_n(q)C\setminus D}
   +\abs{D\setminus\sigma_n(q)C}\\
 &\leq 2e_{q,n}(C)+\bigl|\abs D-\abs C\bigr|\\
 &=O(\eta_n)\abs C.
 \end{aligned}
\]
This proves (i) and (ii).

For \(n\) large, two distinct good source blocks cannot have the same target.  Their images under the bijection \(\sigma_n(q)\) are disjoint, while by \eqref{eq:symmetric-matching} and \eqref{eq:matched-size-ratio} each would occupy a \((1-o_n(1))\)-fraction of the same target block.  Thus \(\Phi_{q,n}\) is injective on \(\cG_{q,n}\).  Finally,
\[
 \sum_{C\in\cG_{q,n}}\abs{\Phi_{q,n}(C)}
 =(1-o_n(1))\sum_{C\in\cG_{q,n}}\abs C
 =(1-o_n(1))N_n,
\]
which proves (iv).
\end{proof}

\subsection{Selection of one component}

We need a standard weighted diagonal selection.  We state it in the form used below.

\begin{lemma}[Weighted component selection]\label{lem:weighted-selection}
For each \(n\), let
\[
 \bigl\{(C_{n,i},D_{n,i}):i\in I_n\bigr\}
\]
be pairs of subsets of a finite set \(Y_n\).  Assume that the \(C_{n,i}\) are pairwise disjoint, the \(D_{n,i}\) are pairwise disjoint, and
\[
 \sum_i\abs{D_{n,i}}=(1-o_n(1))\abs{Y_n}.
\]
For each \(r\in\N\), suppose nonnegative errors \(E_{n,i}^{(r)}\) satisfy
\[
 \sum_iE_{n,i}^{(r)}=o(\abs{Y_n})
\]
for fixed \(r\).  Suppose also that, for every \(M\), the total mass of the \(D_{n,i}\) with \(\abs{D_{n,i}}\leq M\) is \(o(\abs{Y_n})\).  Then there are indices \(i(n)\) such that
\[
 \abs{D_{n,i(n)}}\longrightarrow\infty
\]
and, for every fixed \(r\),
\[
 E_{n,i(n)}^{(r)}=o(\abs{D_{n,i(n)}}).
\]
\end{lemma}

\begin{proof}
Choose \(r_n\to\infty\) so slowly that
\[
 \sum_i\sum_{r\leq r_n}E_{n,i}^{(r)}=o(\abs{Y_n}).
\]
Choose \(\delta_n\downarrow0\) so slowly that the left side is \(o(\delta_n\abs{Y_n})\).  The total mass of indices for which
\[
 \sum_{r\leq r_n}E_{n,i}^{(r)}>\delta_n\abs{D_{n,i}}
\]
is \(o(\abs{Y_n})\).  Slowing \(r_n\) further if necessary, the total mass of the \(D_{n,i}\) with \(\abs{D_{n,i}}\leq r_n\) is also \(o(\abs{Y_n})\); discard these indices as well.  The retained targets still carry mass \((1-o_n(1))\abs{Y_n}\), so some index remains, and any remaining index \(i(n)\) satisfies \(\abs{D_{n,i(n)}}>r_n\to\infty\) and, for each fixed \(r\), eventually \(E^{(r)}_{n,i(n)}\leq\delta_n\abs{D_{n,i(n)}}\).
\end{proof}

We also verify the size hypothesis in our situation.

\begin{lemma}[Component sizes diverge in mass]\label{lem:block-sizes-diverge}
For every fixed \(M\), the union of the \(\cP_n\)-blocks of size at most \(M\) has size \(o(N_n)\).
\end{lemma}

\begin{proof}
Choose a finite set \(E\subseteq\Gamma\) with \(\abs E>M\), possible because \(\Gamma\) is infinite.  For distinct \(g,h\in E\), approximate multiplicativity and faithfulness of \(h^{-1}g\ne1\) give \(\distH(\sigma_n(g),\sigma_n(h))\to1\); taking the finite union over pairs, for all but \(o(N_n)\) points \(y\), the points \(\sigma_n(g)y\), \(g\in E\), are pairwise distinct.  By \eqref{eq:Gamma-block-invariance}, for all but \(o(N_n)\) points they all lie in \(P_n(y)\).  Hence all but \(o(N_n)\) points lie in blocks of size at least \(\abs E>M\).
\end{proof}

Finally, we record how local estimates on an almost invariant set produce a sofic approximation by completion.

\begin{lemma}[Localization by completion]\label{lem:localization-completion}
Let \(H\) be a countable group, let \(\sigma_n:H\to\Sym(Y_n)\) be maps with \(\sigma_n(1)=\mathrm{id}\), and let \(D_n\subseteq Y_n\) be sets with \(\abs{D_n}\to\infty\) such that, for every fixed \(g,h\in H\),
\begin{align*}
 \abs{\set{x\in D_n:\sigma_n(g)x\notin D_n}}&=o(\abs{D_n}),\\
 \abs{\set{x\in D_n:\sigma_n(g)\sigma_n(h)x\ne\sigma_n(gh)x}}&=o(\abs{D_n}),\\
 \abs{\set{x\in D_n:\sigma_n(g)x=x}}&=o(\abs{D_n})\quad(g\ne1).
\end{align*}
For \(g\in H\), the restriction of \(\sigma_n(g)\) to \(D_n\cap\sigma_n(g)^{-1}D_n\) is a partial injection of \(D_n\); complete it to a permutation \(\pi_n(g)\in\Sym(D_n)\).  Then \(\pi_n:H\to\Sym(D_n)\) is a sofic approximation of \(H\).  Moreover, on any finite symmetric set \(E\subseteq H\) of elements for which \(\sigma_n(g^{-1})=\sigma_n(g)^{-1}\), the completions may be chosen so that \(\pi_n(g^{-1})=\pi_n(g)^{-1}\) for all \(g\in E\).
\end{lemma}

\begin{proof}
Fix \(g\in H\) and let \(A_n(g)=\set{x\in D_n:\pi_n(g)x\ne\sigma_n(g)x}\).  Since \(\pi_n(g)\) agrees with \(\sigma_n(g)\) on \(D_n\cap\sigma_n(g)^{-1}D_n\), the set \(A_n(g)\) is contained in the escape set, so \(\abs{A_n(g)}=o(\abs{D_n})\).

\emph{Multiplicativity.}  For fixed \(g,h\in H\) and \(x\in D_n\) outside
\[
 A_n(h)\cup A_n(gh)\cup \sigma_n(h)^{-1}A_n(g)\cup
 \set{x:\sigma_n(g)\sigma_n(h)x\ne\sigma_n(gh)x},
\]
we have \(\pi_n(h)x=\sigma_n(h)x\), then \(\pi_n(g)\pi_n(h)x=\sigma_n(g)\sigma_n(h)x=\sigma_n(gh)x=\pi_n(gh)x\).  The excluded set has size \(o(\abs{D_n})\); for the third term, note that \(\sigma_n(h)^{-1}A_n(g)\cap D_n\) is covered by the escape set of \(h\), the set \(A_n(h)\), and the preimage under the partial bijection of \(A_n(g)\), each of size \(o(\abs{D_n})\).

\emph{Freeness.}  For \(g\ne1\), the fixed points of \(\pi_n(g)\) are contained in \(A_n(g)\cup\set{x\in D_n:\sigma_n(g)x=x}\), of size \(o(\abs{D_n})\).

\emph{Inverse compatibility.}  Let \(g\in E\).  If \(g\ne g^{-1}\), complete \(g\) arbitrarily and define \(\pi_n(g^{-1})=\pi_n(g)^{-1}\); since \(\sigma_n(g^{-1})=\sigma_n(g)^{-1}\), the permutation \(\pi_n(g^{-1})\) agrees with \(\sigma_n(g^{-1})\) on \(\sigma_n(g)(D_n\cap\sigma_n(g)^{-1}D_n)=D_n\cap\sigma_n(g^{-1})^{-1}D_n\), as required.  If \(g=g^{-1}\), then \(\sigma_n(g)\) is an involution, its partial restriction is an involution of the \(\sigma_n(g)\)-invariant set \(D_n\cap\sigma_n(g)^{-1}D_n\), and extending by any involution of the complement (for instance the identity) makes \(\pi_n(g)\) an involution.
\end{proof}

\subsection{Proof of the criterion}

\begin{proof}[Proof of Theorem~\ref{thm:criterion}]
Apply Proposition~\ref{prop:component-matching} with \(q=q_0\).  Write the good matched pairs as
\[
 (C,D),
 \qquad
 D=\Phi_{q_0,n}(C).
\]
The source blocks are disjoint, the target blocks are disjoint, and the targets have total mass \((1-o_n(1))N_n\).

Let
\[
 L_1\subseteq L_2\subseteq\cdots\subseteq KJ
\]
be a finite symmetric exhaustion containing \(1\), arranged so that products needed to test multiplication on \(L_r\) lie in \(L_{r+1}\).  For each matched pair and each fixed \(r\), define an error \(E_{n,C,D}^{(r)}\) to be the sum of:
\begin{enumerate}[label=\textup{(\alph*)}]
\item the matching error \(\abs{\sigma_n(q_0)C\symdiff D}\);
\item the number of edges in the symmetric difference of \(\widetilde X_n\) and the \(S\)-generator graph that are incident to \(C\);
\item the conjugacy defects in \eqref{eq:conjugacy-defect}, restricted to \(C\), for \(s\in S\);
\item the multiplication defects for pairs in \(L_r\), restricted to \(D\);
\item for \(1\ne g\in L_r\), the number of points of \(D\) fixed by \(\sigma_n(g)\);
\item for \(g\in L_r\), the number of points of \(D\) sent outside \(D\) by \(\sigma_n(g)\).
\end{enumerate}
For fixed \(r\), the sum of these errors over all matched pairs is \(o(N_n)\): the source blocks are disjoint, so (b) is bounded by twice the global edge modification and (c) by the global conjugacy defects \eqref{eq:conjugacy-defect}; the target blocks are disjoint, so global soficity controls (d) and (e); and \eqref{eq:Gamma-block-invariance} controls (f), because
\[
 KJ\leq\Gamma
\]
and every target is a \(\cP_n\)-block.  The matching estimate handles (a).  Lemmas~\ref{lem:weighted-selection} and \ref{lem:block-sizes-diverge} therefore give matched pairs
\begin{equation}\label{eq:selected-pairs}
 (C_n,D_n),
 \qquad
 D_n=\Phi_{q_0,n}(C_n),
\end{equation}
such that
\begin{equation}\label{eq:selected-size}
 \abs{D_n}\to\infty
\end{equation}
and every fixed finite family of the errors above is \(o(\abs{D_n})\).

Put
\[
 U_n=\sigma_n(q_0)C_n.
\]
We know
\begin{equation}\label{eq:U-D-close}
 \abs{U_n\symdiff D_n}=o(\abs{D_n}).
\end{equation}
The sets need not have exactly the same cardinality.  Choose \(D'_n\subseteq Y_n\) such that
\begin{equation}\label{eq:D-prime}
 \abs{D'_n}=\abs{U_n}
 \quad\text{and}\quad
 \abs{D'_n\symdiff D_n}\leq\abs{U_n\symdiff D_n}.
\end{equation}
If \(\abs{D_n}>\abs{U_n}\), remove the required number of points from \(D_n\setminus U_n\); if \(\abs{D_n}<\abs{U_n}\), add points from \(U_n\setminus D_n\).  Every selected local estimate transfers from \(D_n\) to \(D'_n\), since changing a set by \(r\) points changes each fixed permutation escape count, fixed-point count, or multiplication-defect count by at most \(O(r)\).

By the transfer estimates just noted, the hypotheses of Lemma~\ref{lem:localization-completion} hold for \(H=KJ\) and the sets \(D'_n\); note that \(\abs{D'_n}=\abs{U_n}=\abs{C_n}\to\infty\).  Choosing the completions compatibly on the finite symmetric set \(S_K\), as permitted by that lemma, we obtain a sofic approximation
\[
 \pi_n:KJ\longrightarrow\Sym(D'_n).
\]
Since \([K,J]=1\) and \(K\cap J=1\), multiplication induces an isomorphism
\[
 KJ\cong K\times J.
\]

It remains to show that the \(K\)-generator graph is essentially one expander.  Transport the expander \(\widetilde X_n[C_n]\) by \(\sigma_n(q_0)\), obtaining an \(h_\Gamma\)-expander on \(U_n\).  Because \(\abs{U_n}=\abs{D'_n}\), choose a bijection
\[
 \theta_n:U_n\longrightarrow D'_n
\]
which fixes \(U_n\cap D'_n\).  The number of moved vertices is \(o(\abs{D'_n})\) by \eqref{eq:U-D-close} and \eqref{eq:D-prime}.  Push the transported expander through \(\theta_n\).  This remains an \(h_\Gamma\)-expander on \(D'_n\).

Let \(H_n\) denote the pushed graph and \(X_{K,n}\) the \(S_K\)-generator graph of \(\pi_n\).  Their edge multisets can disagree for exactly four reasons:
\begin{enumerate}[label=\textup{(\arabic*)}]
\item source edges incident to \(C_n\) that were changed in passing from the \(S\)-generator graph to \(\widetilde X_n\);
\item the conjugacy defects \(\sigma_n(q_0)\sigma_n(s)y\ne\sigma_n(q_0sq_0^{-1})\sigma_n(q_0)y\) at points \(y\in C_n\), for \(s\in S\);
\item edges incident to vertices moved by \(\theta_n\);
\item points at which a completed permutation \(\pi_n(k)\), \(k\in S_K\), differs from \(\sigma_n(k)\).
\end{enumerate}
Items (1) and (2) are part of the selected error budget; item (3) is \(o(\abs{D'_n})\) because the pushed graph has maximum degree at most \(d_\Gamma\) and \(\theta_n\) moves \(o(\abs{D'_n})\) vertices; and item (4) is \(o(\abs{D'_n})\) by Lemma~\ref{lem:localization-completion} and the transfer estimates.  Hence
\[
 \abs{E(H_n)\symdiff E(X_{K,n})}=o(\abs{D'_n}).
\]
Thus the \(S_K\)-edges of \(\pi_n\) form an essential expander sequence.  The group \(K\cong\Gamma\) has property~\((T)\), and \(J\) is finitely generated.  Theorem~\ref{thm:Kun-Thom}, applied to the sofic approximation of \(K\times J\) just constructed, implies that \(J\) is LEF.
\end{proof}

\begin{remark}\label{rem:product-obstruction}
The argument never claims that the ambient generators fix every \(\Gamma\)-component.  They may permute equal-sized components, as in a product approximation \(X_n\times Z_n\).  The proof instead chooses one matched target component which is almost preserved by \(KJ\leq\Gamma\), while the transported source component supplies expansion there.  This is the point that avoids the standard product obstruction.
\end{remark}

\begin{remark}[Role of the hypothesis \(J\leq\Gamma\)]\label{rem:J-in-Gamma-role}
The hypothesis \(J\leq\Gamma\) enters the proof exactly once, through \eqref{eq:Gamma-block-invariance}: it is what confines the \(J\)-action to the selected target component.  It is not a technical convenience.  In a product approximation of \(\Gamma\times\Delta\), the \(\Delta\)-side permutes the \(\Gamma\)-components, and a finitely generated group commuting with \(\Gamma\) need not be LEF when it acts transversally to the component structure: the Baumslag--Solitar group \(\mathrm{BS}(2,3)\) is sofic, being residually solvable, but it is finitely presented and non-Hopfian \cite{BaumslagSolitar62}, hence not residually finite and, by Lemma~\ref{lem:fp-lef-rf}, not LEF.
\end{remark}

\section{The binary Leavitt algebra and its elementary groups}\label{sec:leavitt}

Let \(R=L_{\F_2}(1,2)\) be the binary Leavitt algebra \cite{Leavitt62}, given by \eqref{eq:Leavitt-presentation}.  Put
\[
 p_0=s_0t_0,
 \qquad
 p_1=s_1t_1.
\]
Then \(p_0,p_1\) are orthogonal idempotents with \(p_0+p_1=1\).

For a word \(\alpha=s_{i_1}\cdots s_{i_r}\), set
\[
 \alpha^*=t_{i_r}\cdots t_{i_1}.
\]
If \(\alpha\) and \(\beta\) are prefix-incomparable, then \(\alpha^*\beta=0\), while \(\alpha^*\alpha=1\).

\subsection{Elementary two-by-two identities}\label{subsec:two-by-two}

All invertible matrices in this paper are shown to be products of elementary matrices by explicit factorization; no \(K\)-theory is used.  Since \(R\) is an \(\F_2\)-algebra, signs may be ignored throughout.  We record three identities in \(M_2(R)\); they will be embedded into a pair of coordinates of \(M_4(R)\) as needed.

\begin{lemma}\label{lem:two-by-two}
In \(M_2(R)\), write \(x_{12}(a)=\begin{psmallmatrix}1&a\\0&1\end{psmallmatrix}\) and \(x_{21}(a)=\begin{psmallmatrix}1&0\\a&1\end{psmallmatrix}\).  Then:
\begin{enumerate}[label=\textup{(\alph*)}]
\item \(x_{12}(s_1)\,x_{21}(t_1)\,x_{12}(s_1)=\begin{pmatrix}p_0&s_1\\t_1&0\end{pmatrix}\);
\item \(x_{21}(1+t_0)\,x_{12}(1)\,x_{21}(1+s_0)\,x_{12}(t_0)=\begin{pmatrix}s_0&s_1t_1\\0&t_0\end{pmatrix}\);
\item for every unit \(a\in R^\times\),
\[
 w(a):=x_{12}(a)\,x_{21}(a^{-1})\,x_{12}(a)=\begin{pmatrix}0&a\\a^{-1}&0\end{pmatrix},
 \qquad
 \diag(a,a^{-1})=w(a)\,w(1)\in\EL_2(R),
\]
and consequently, for all units \(a,b\in R^\times\),
\[
 \diag\bigl(aba^{-1}b^{-1},1\bigr)
 =\diag(a,a^{-1})\,\diag(b,b^{-1})\,\diag\bigl((ba)^{-1},ba\bigr)\in\EL_2(R).
\]
\end{enumerate}
\end{lemma}

\begin{proof}
Everything follows by direct computation from \(t_is_j=\delta_{ij}\), \(p_0+p_1=1\), and characteristic \(2\).

(a) First, \(x_{12}(s_1)x_{21}(t_1)=\begin{psmallmatrix}1+s_1t_1&s_1\\t_1&1\end{psmallmatrix}=\begin{psmallmatrix}p_0&s_1\\t_1&1\end{psmallmatrix}\).  Right multiplication by \(x_{12}(s_1)\) adds the first column times \(s_1\) to the second; since \(p_0s_1=s_0t_0s_1=0\) and \(t_1s_1=1\), the second column becomes \(\begin{psmallmatrix}s_1\\0\end{psmallmatrix}\).

(b) First, \(x_{21}(1+t_0)x_{12}(1)=\begin{psmallmatrix}1&1\\1+t_0&t_0\end{psmallmatrix}\).  Right multiplication by \(x_{21}(1+s_0)\) adds the second column times \(1+s_0\) to the first; using \(t_0s_0=1\), the first column becomes \(\begin{psmallmatrix}1+(1+s_0)\\(1+t_0)+(t_0+1)\end{psmallmatrix}=\begin{psmallmatrix}s_0\\0\end{psmallmatrix}\).  Right multiplication by \(x_{12}(t_0)\) then adds the first column times \(t_0\) to the second, which becomes \(\begin{psmallmatrix}1+s_0t_0\\t_0\end{psmallmatrix}=\begin{psmallmatrix}s_1t_1\\t_0\end{psmallmatrix}\).

(c) First, \(x_{12}(a)x_{21}(a^{-1})=\begin{psmallmatrix}1+1&a\\a^{-1}&1\end{psmallmatrix}=\begin{psmallmatrix}0&a\\a^{-1}&1\end{psmallmatrix}\), and right multiplication by \(x_{12}(a)\) turns the second column into \(\begin{psmallmatrix}0\cdot a+a\\a^{-1}a+1\end{psmallmatrix}=\begin{psmallmatrix}a\\0\end{psmallmatrix}\).  The identity \(w(a)w(1)=\diag(a,a^{-1})\) and the final display are multiplications of diagonal and antidiagonal matrices.
\end{proof}

The ring \(R\) is finitely generated as an associative unital ring.  We shall use the following theorem.

\begin{theorem}[Ershov--Jaikin-Zapirain \cite{ErshovJaikin10}]\label{thm:EJZ}
If \(A\) is a finitely generated associative unital ring and \(m\geq3\), then \(\EL_m(A)\) has property~\((T)\).
\end{theorem}

In \cite{ErshovJaikin10} this is proved for the universal rings \(\Z\langle x_1,\dots,x_k\rangle\); the general case follows because \(\EL_m(A)\) is a quotient of \(\EL_m(\Z\langle x_1,\dots,x_k\rangle)\) whenever \(A\) is generated by \(k\) elements as a unital ring, and property \((T)\) passes to quotients.

Set
\begin{equation}\label{eq:def-G}
 G=\EL_4(R)
\end{equation}
and let
\begin{equation}\label{eq:def-Gamma}
 \Gamma=
 \set{
 \begin{pmatrix}\gamma&0\\0&1\end{pmatrix}
 \suchthat \gamma\in\EL_3(R)}.
\end{equation}
By Theorem~\ref{thm:EJZ}, both \(G\) and \(\Gamma\cong\EL_3(R)\) have property~\((T)\).

\section{Two compressors with a common image}\label{sec:compressors}

\subsection{A four-leaf binary partition}

Define
\begin{equation}\label{eq:alpha-leaves}
 \alpha_i=s_0^{i-1}s_1\quad(1\leq i\leq3),
 \qquad
 \alpha_4=s_0^3.
\end{equation}
These are the leaves obtained by repeatedly expanding the leftmost leaf in the binary rooted tree.  They satisfy
\begin{equation}\label{eq:leaf-orthogonality}
 \alpha_i^*\alpha_j=\delta_{ij}
\end{equation}
and
\begin{equation}\label{eq:leaf-sum}
 \sum_{i=1}^4\alpha_i\alpha_i^*=1.
\end{equation}
For \eqref{eq:leaf-sum}, telescope:
\[
 \begin{aligned}
 \sum_{i=1}^3s_0^{i-1}s_1t_1t_0^{i-1}+s_0^3t_0^3
 &=\sum_{j=0}^2s_0^j(1-s_0t_0)t_0^j+s_0^3t_0^3\\
 &=\sum_{j=0}^2\bigl(s_0^jt_0^j-s_0^{j+1}t_0^{j+1}\bigr)+s_0^3t_0^3\\
 &=1.
 \end{aligned}
\]

\subsection{The first compressor}

Let
\[
 \mathbf b=
 \begin{pmatrix}
 s_1\alpha_1^*\\ \vdots\\ s_1\alpha_3^*
 \end{pmatrix},
 \qquad
 \mathbf c=
 \begin{pmatrix}
 \alpha_1t_1&\cdots&\alpha_3t_1
 \end{pmatrix}.
\]
Define
\begin{equation}\label{eq:def-u}
 u=
 \begin{pmatrix}
 s_0I_3&\mathbf b\\
 0&\alpha_4^*
 \end{pmatrix}.
\end{equation}

\begin{proposition}[Explicit elementary factorization of \(u\)]\label{prop:u-elementary}
For \(1\leq i\leq3\), put
\[
 U_i=x_{4i}(1+t_0)\,x_{i4}(1)\,x_{4i}(1+s_0)\,x_{i4}(t_0)\in\EL_4(R).
\]
Then
\[
 u=U_3U_2U_1.
\]
In particular, \(u\in G\).
\end{proposition}

\begin{proof}
By Lemma~\ref{lem:two-by-two}(b), applied in the pair of coordinates \((i,4)\), the matrix \(U_i\) is the identity outside those coordinates and equals \(\begin{psmallmatrix}s_0&s_1t_1\\0&t_0\end{psmallmatrix}\) on them; explicitly,
\[
 U_i=I_4+(s_0+1)E_{ii}+s_1t_1E_{i4}+(t_0+1)E_{44},
\]
the signs being immaterial in characteristic \(2\).  Put \(P_k=U_kU_{k-1}\cdots U_1\).  We prove by induction on \(k\) that
\[
 P_k=\sum_{i\leq k}\bigl(s_0E_{ii}+s_1t_1t_0^{\,i-1}E_{i4}\bigr)
 +\sum_{k<i\leq3}E_{ii}
 +t_0^{\,k}E_{44}.
\]
The case \(k=1\) is the displayed form of \(U_1\).  For the inductive step, compute the rows of \(P_{k+1}=U_{k+1}P_k\).  For \(i\notin\{k+1,4\}\), row \(i\) of \(U_{k+1}\) is \(e_i^{\mathsf T}\), so row \(i\) of \(P_{k+1}\) equals row \(i\) of \(P_k\).  Row \(k+1\) of \(U_{k+1}\) is \(s_0e_{k+1}^{\mathsf T}+s_1t_1e_4^{\mathsf T}\); since row \(k+1\) of \(P_k\) is \(e_{k+1}^{\mathsf T}\) and row \(4\) of \(P_k\) is \(t_0^{\,k}e_4^{\mathsf T}\), row \(k+1\) of \(P_{k+1}\) is \(s_0e_{k+1}^{\mathsf T}+s_1t_1t_0^{\,k}e_4^{\mathsf T}\).  Row \(4\) of \(U_{k+1}\) is \(t_0e_4^{\mathsf T}\), so row \(4\) of \(P_{k+1}\) is \(t_0^{\,k+1}e_4^{\mathsf T}\).  This is the claimed formula with \(k+1\) in place of \(k\).

Since \(s_1t_1t_0^{\,i-1}=s_1\alpha_i^*\) for \(i\leq3\) and \(t_0^{\,3}=\alpha_4^*\), the matrix \(P_3\) is exactly \(u\) as defined in \eqref{eq:def-u}.
\end{proof}

\begin{lemma}\label{lem:u-inverse}
The matrix \(u\) is invertible, with
\begin{equation}\label{eq:u-inverse}
 u^{-1}=
 \begin{pmatrix}
 t_0I_3&0\\
 \mathbf c&\alpha_4
 \end{pmatrix}.
\end{equation}
\end{lemma}

\begin{proof}
Using \eqref{eq:leaf-orthogonality},
\[
 (\mathbf b\mathbf c)_{ij}
 =s_1\alpha_i^*\alpha_jt_1
 =\delta_{ij}p_1.
\]
Moreover,
\[
 \mathbf b\alpha_4=0,
 \qquad
 \alpha_4^*\mathbf c=0,
 \qquad
 t_0\mathbf b=0,
 \qquad
 \mathbf c s_0=0,
\]
and
\[
 \mathbf c\mathbf b
 =\sum_{i=1}^3\alpha_i\alpha_i^*
 =1-\alpha_4\alpha_4^*.
\]
The four vanishing identities follow from prefix orthogonality and \(t_0s_1=t_1s_0=0\).  Together with \(s_0t_0+p_1=1\), \(t_0s_0=1\), and \eqref{eq:leaf-sum}, these identities show that the matrices in \eqref{eq:def-u} and \eqref{eq:u-inverse} multiply to \(I_4\) in both orders.
\end{proof}

For \(\gamma=(\gamma_{ij})\in M_3(R)\), write \(s_0\gamma t_0\) for the matrix with entries \(s_0\gamma_{ij}t_0\).

\begin{proposition}\label{prop:u-compresses}
For \(\gamma\in\EL_3(R)\),
\begin{equation}\label{eq:u-conjugation-general}
 u
 \begin{pmatrix}\gamma&0\\0&1\end{pmatrix}
 u^{-1}
 =
 \begin{pmatrix}
 p_1I_3+s_0\gamma t_0&0\\
 0&1
 \end{pmatrix}.
\end{equation}
Consequently,
\begin{equation}\label{eq:def-K}
 K:=u\Gamma u^{-1}
 =\angles{x_{ij}(s_0at_0)
 \suchthat a\in R,\ 1\leq i\ne j\leq3}
 \leq\Gamma.
\end{equation}
\end{proposition}

\begin{proof}
Multiplying the block matrices gives top-left block
\[
 s_0\gamma t_0+\mathbf b\mathbf c
 =s_0\gamma t_0+p_1I_3.
\]
The off-diagonal blocks vanish by prefix orthogonality, and the bottom-right entry is \(1\).  For an elementary generator \(\gamma=I_3+aE_{ij}\), the top-left block becomes
\[
 p_1I_3+s_0I_3t_0+s_0a t_0E_{ij}
 =I_3+s_0at_0E_{ij}.
\]
This proves the description of \(K\), which lies in the upper-left \(\EL_3(R)\)-copy.
\end{proof}

\subsection{The second compressor and generation of \texorpdfstring{\(G\)}{G}}

Let \(z\in M_4(R)\) act as the identity on coordinates \(2\) and \(3\), and let its block on coordinates \(1,4\) be
\begin{equation}\label{eq:z-block}
 \begin{pmatrix}
 p_0&s_1\\
 t_1&0
 \end{pmatrix}.
\end{equation}

\begin{lemma}\label{lem:z-properties}
The matrix \(z\) satisfies
\[
 z=x_{14}(s_1)\,x_{41}(t_1)\,x_{14}(s_1)\in G;
\]
it is an involution and it centralizes \(K\).
\end{lemma}

\begin{proof}
The factorization is Lemma~\ref{lem:two-by-two}(a), applied in the pair of coordinates \((1,4)\); the block of \(z\) on those coordinates is \eqref{eq:z-block} and \(z\) fixes the other coordinates.  The square of the block \eqref{eq:z-block} is
\[
 \begin{pmatrix}
 p_0^2+s_1t_1&p_0s_1\\
 t_1p_0&t_1s_1
 \end{pmatrix}
 =
 \begin{pmatrix}1&0\\0&1\end{pmatrix}.
\]
Thus \(z^2=1\).

It is enough to check the elementary generators of \(K\).  Put \(r=s_0at_0\).  Then
\[
 p_0r=r,
 \quad
 t_1r=0,
 \quad
 rp_0=r,
 \quad
 rs_1=0.
\]
Direct multiplication gives
\[
 zx_{1j}(r)z=x_{1j}(r),
 \qquad
 zx_{j1}(r)z=x_{j1}(r)
 \quad(2\leq j\leq3),
\]
and generators not involving the first coordinate are fixed trivially.
\end{proof}

Set
\begin{equation}\label{eq:def-v}
 v=zu.
\end{equation}
Then
\begin{equation}\label{eq:common-image}
 v\Gamma v^{-1}=zKz^{-1}=K=u\Gamma u^{-1}.
\end{equation}

\begin{proposition}\label{prop:generation}
The ambient group is generated by the subgroup and the two compressors:
\begin{equation}\label{eq:G-generated}
 G=\langle\Gamma,u,v\rangle.
\end{equation}
\end{proposition}

\begin{proof}
Since \(z=vu^{-1}\), it suffices to prove \(\langle\Gamma,z\rangle=G\).  For \(2\leq j\leq3\) and \(a\in R\), direct multiplication with the block \eqref{eq:z-block} gives
\begin{equation}\label{eq:z-row-generation}
 zx_{1j}(s_1a)z=x_{4j}(a)
\end{equation}
and
\begin{equation}\label{eq:z-column-generation}
 zx_{j1}(at_1)z=x_{j4}(a).
\end{equation}
Thus \(\langle\Gamma,z\rangle\) contains every elementary matrix \(x_{4j}(a)\) and \(x_{j4}(a)\) with \(2\leq j\leq3\).  The Steinberg commutator identities yield
\[
 [x_{12}(a),x_{24}(1)]=x_{14}(a),
 \qquad
 [x_{42}(a),x_{21}(1)]=x_{41}(a).
\]
Hence all elementary matrices involving the fourth coordinate lie in \(\langle\Gamma,z\rangle\), as do all elementary matrices on the first three coordinates.  These generate \(\EL_4(R)\).
\end{proof}

\section{A commuting copy of Thompson's group \texorpdfstring{\(V\)}{V}}\label{sec:V}

A \emph{leaf set} is a finite set \(\Lambda\) of binary words such that every infinite binary word has exactly one prefix in \(\Lambda\); equivalently, \(\Lambda\) is the set of leaves of a finite complete rooted binary tree.  For every leaf set \(\Lambda\),
\begin{equation}\label{eq:leafset-identities}
 \sum_{\lambda\in\Lambda}\lambda\lambda^*=1,
 \qquad
 \lambda^*\mu=\delta_{\lambda\mu}\quad(\lambda,\mu\in\Lambda),
\end{equation}
the first by induction on the tree using \(s_0t_0+s_1t_1=1\), the second by prefix orthogonality.

An element \(g\) of Thompson's group \(V\) is given by a \emph{table}: a bijection \(\beta_i\mapsto\alpha_i\) between two leaf sets, acting on infinite binary words by the prefix substitutions \(\beta_i\xi\mapsto\alpha_i\xi\).  Refining a table replaces a pair \((\alpha_i,\beta_i)\) by the two pairs \((\alpha_i0,\beta_i0)\) and \((\alpha_i1,\beta_i1)\); two tables define the same element of \(V\) if and only if they admit a common refinement \cite{CannonFloydParry96}.  Define
\begin{equation}\label{eq:V-in-R}
 g\longmapsto U_g=\sum_i\alpha_i\beta_i^*\in R.
\end{equation}

\begin{proposition}\label{prop:V-embeds}
The assignment \eqref{eq:V-in-R} is a well-defined injective group homomorphism from \(V\) into \(R^\times\).
\end{proposition}

\begin{proof}
\emph{Well defined.}  Refining one pair changes the sum by \(\alpha_i(s_0t_0+s_1t_1)\beta_i^*-\alpha_i\beta_i^*=0\), so \(U_g\) is invariant under refinement, and any two tables for \(g\) yield the same element of \(R\).

\emph{Multiplicative.}  Let \(g\) have table \(\beta_i\mapsto\alpha_i\) and \(h\) have table \(\delta_j\mapsto\gamma_j\).  Refining, we may assume \(\{\gamma_j\}=\{\beta_i\}\) as leaf sets, say \(\gamma_{\tau(i)}=\beta_i\) for a bijection \(\tau\).  By \eqref{eq:leafset-identities},
\[
 U_gU_h=\sum_{i,j}\alpha_i\,(\beta_i^*\gamma_j)\,\delta_j^*
 =\sum_i\alpha_i\delta_{\tau(i)}^*,
\]
which is \(U_{gh}\) for the composition \(\xi\mapsto g(h(\xi))\).  Moreover \(U_1=\sum_{\lambda\in\Lambda}\lambda\lambda^*=1\) for any leaf set \(\Lambda\), so each \(U_g\) is a unit with \(U_g^{-1}=U_{g^{-1}}\).

\emph{Injective.}  Let \(R\) act on the \(\F_2\)-vector space with basis the set of infinite binary words: \(s_i\) prefixes the letter \(i\), while \(t_i\) deletes an initial letter \(i\) and annihilates the words beginning with the other letter.  The defining relations of \eqref{eq:Leavitt-presentation} hold in this action, so it is a representation of \(R\), and by \eqref{eq:leafset-identities} the element \(U_g\) sends each basis word \(\beta_j\xi\) to \(\alpha_j\xi\).  If \(U_g=1\), then the action fixes every infinite binary word, so the prefix substitution defined by the table of \(g\) is the identity homeomorphism of the Cantor set, and \(g=1\) in \(V\).  Hence the kernel is trivial and the homomorphism is injective.
\end{proof}

This is the standard Leavitt-algebra realization of the Higman--Thompson groups; see also \cite{Nekrashevych04,Pardo11}.

From now on, identify each \(g\in V\) with the unit \(U_g\in R^\times\).  For \(g\in V\leq R^\times\), define
\begin{equation}\label{eq:def-cg}
 c(g)=p_0+s_1gt_1.
\end{equation}

\begin{lemma}\label{lem:corner-V}
The map \(c:V\to R^\times\) is an injective homomorphism.  Moreover,
\begin{equation}\label{eq:c-fixes-corner}
 c(g)s_0=s_0,
 \qquad
 t_0c(g)=t_0.
\end{equation}
\end{lemma}

\begin{proof}
Orthogonality of the two corners gives
\[
 c(g)c(h)=p_0+s_1ght_1=c(gh)
\]
and \(c(g)^{-1}=c(g^{-1})\).  Also
\[
 t_1c(g)s_1=g,
\]
so the map is injective.  Finally,
\[
 c(g)s_0=p_0s_0+s_1gt_1s_0=s_0,
\]
and similarly \(t_0c(g)=t_0\).
\end{proof}

Define
\begin{equation}\label{eq:def-J}
 j(g)=\diag(c(g),1,1,1)\in\GL_4(R),
\end{equation}
where the final coordinate is the fourth coordinate and the upper-left \(3\times3\) block is \(\diag(c(g),1,1)\).

\begin{proposition}\label{prop:J-in-Gamma}
For every \(g\in V\),
\[
 \diag(c(g),1,1)\in\EL_3(R),
\]
so
\begin{equation}\label{eq:J-in-Gamma}
 J:=j(V)\leq\Gamma .
\end{equation}
\end{proposition}

\begin{proof}
The group \(V\) is simple \cite{Higman74,CannonFloydParry96} and infinite, hence nonabelian, hence perfect: \(V=[V,V]\).  Write \(g=\prod_{i=1}^{\ell}[a_i,b_i]\) with \(a_i,b_i\in V\).  Since \(c\) is a homomorphism by Lemma~\ref{lem:corner-V},
\[
 c(g)=\prod_{i=1}^{\ell}\bigl[c(a_i),c(b_i)\bigr]
\]
is a product of commutators of units of \(R\).  By Lemma~\ref{lem:two-by-two}(c), \(\diag([x,y],1)\in\EL_2(R)\) for all units \(x,y\), so \(\diag(c(g),1)\in\EL_2(R)\).  Embedding \(\EL_2(R)\) into the coordinates \((1,2)\) of \(\EL_3(R)\) gives the claim.
\end{proof}

\begin{proposition}\label{prop:J-centralizes-K}
The subgroups \(K\) and \(J\) satisfy
\[
 [K,J]=1
 \qquad\text{and}\qquad
 K\cap J=\{1\}.
\]
Consequently,
\[
 KJ\cong K\times J.
\]
\end{proposition}

\begin{proof}
Let \(r=s_0at_0\).  Equation \eqref{eq:c-fixes-corner}, applied also to \(g^{-1}\), gives
\[
 c(g)r=r,
 \qquad
 rc(g)^{-1}=r.
\]
Thus conjugation by \(j(g)\) fixes every generator \(x_{ij}(r)\) of \(K\), proving \([K,J]=1\).

By \eqref{eq:u-conjugation-general}, every \(k\in K\) has upper-left block
\[
 p_1I_3+s_0\gamma t_0
\]
for some \(\gamma\in\EL_3(R)\).  Hence
\begin{equation}\label{eq:K-sandwich-zero}
 t_1(k_{11}-1)s_1=0.
\end{equation}
On the other hand,
\begin{equation}\label{eq:J-sandwich}
 t_1(c(g)-1)s_1=g-1.
\end{equation}
If \(j(g)\in K\), equations \eqref{eq:K-sandwich-zero} and \eqref{eq:J-sandwich} imply \(g=1\).  Thus the intersection is trivial.
\end{proof}

\begin{proposition}\label{prop:V-not-LEF}
The group \(J\cong V\) is not LEF.
\end{proposition}

\begin{proof}
Thompson's group \(V\) is infinite, simple, and finitely presented \cite{Higman74,BleakQuick17}.  By Lemma~\ref{lem:fp-lef-rf}, a finitely presented LEF group is residually finite.  An infinite simple group has no nontrivial finite quotient: the kernel of such a quotient would be a proper normal subgroup, while an injective map into a finite group is impossible.  Hence \(V\), and therefore \(J\), is not LEF.
\end{proof}

\section{Nonsoficity of the elementary group}\label{sec:nonsofic}

\begin{theorem}\label{thm:EL4-nonsofic}
The group
\[
 \EL_4\bigl(L_{\F_2}(1,2)\bigr)
\]
is not sofic.
\end{theorem}

\begin{proof}
Let \(G\), \(\Gamma\), \(u\), \(v\), \(K\), and \(J\) be as above.  The groups \(G\) and \(\Gamma\) have property~\((T)\) by Theorem~\ref{thm:EJZ}.  Equations \eqref{eq:common-image} and \eqref{eq:G-generated} give
\[
 G=\langle\Gamma,u,v\rangle,
 \qquad
 u\Gamma u^{-1}=v\Gamma v^{-1}=K\leq\Gamma.
\]
Proposition~\ref{prop:J-centralizes-K} gives
\[
 J\leq\Gamma,
 \qquad
 [K,J]=1,
 \qquad
 K\cap J=1.
\]
Thus every hypothesis of Theorem~\ref{thm:criterion} holds with \(Q=\{u,v\}\) and \(q_0=u\).  If \(G\) were sofic, that theorem would imply that \(J\) is LEF, contradicting Proposition~\ref{prop:V-not-LEF}.
\end{proof}

\begin{corollary}\label{cor:G-properties}
The group \(G=\EL_4(R)\) is infinite, finitely generated, has property~\((T)\), and is nonsofic.
\end{corollary}

\begin{proof}
Property~\((T)\) was established above.  For a direct finite-generation argument, choose a finite unital ring generating set \(T\) for \(R\).  The identities
\[
 x_{ij}(a+b)=x_{ij}(a)x_{ij}(b),
 \qquad
 [x_{ik}(a),x_{kj}(b)]=x_{ij}(ab)
 \quad(i,j,k\ \text{distinct})
\]
show inductively that the finitely many matrices \(x_{ij}(t)\) with \(t\in T\cup\{1\}\) generate every elementary matrix, so \(G\) is finitely generated.  The group is infinite because it contains \(J\cong V\).  Nonsoficity is Theorem~\ref{thm:EL4-nonsofic}.
\end{proof}

\section{A finitely presented nonsofic group}\label{sec:fp}

We finish with a general finite-table construction.

For a finite set \(F\subseteq G\), a map
\[
 p:F\cup F^2\longrightarrow\Sym(Y)
\]
is an \((F,\eps)\)-model if
\begin{align}
 \distH\bigl(p(g)p(h),p(gh)\bigr)&<\eps
 &&(g,h\in F),\label{eq:model-mult}\\
 \distH\bigl(p(g),p(h)\bigr)&>1-\eps
 &&(g,h\in F,\ g\ne h).\label{eq:model-sep}
\end{align}
\begin{lemma}\label{lem:models-equivalence}
A countable group \(G\) is sofic if and only if, for every finite \(F\subseteq G\) with \(1\in F\) and every \(\eps>0\), there is an \((F,\eps)\)-model.
\end{lemma}

\begin{proof}
Suppose \(G\) is sofic, and fix \(F\) and \(\eps\).  Take a sofic approximation, normalized as in Lemma~\ref{lem:finite-normalization} on the symmetrization of \(F\cup F^2\), and put \(p=\sigma_n|_{F\cup F^2}\) for a large \(n\).  Condition \eqref{eq:model-mult} is approximate multiplicativity.  For \(g\ne h\) in \(F\), bi-invariance of \(\distH\) gives
\[
 \distH\bigl(p(g),p(h)\bigr)
 =\distH\bigl(\id,\sigma_n(g)^{-1}\sigma_n(h)\bigr)
 \geq
 \distH\bigl(\id,\sigma_n(g^{-1}h)\bigr)
 -\distH\bigl(\sigma_n(g^{-1})\sigma_n(h),\sigma_n(g^{-1}h)\bigr),
\]
and the right-hand side tends to \(1\), which gives \eqref{eq:model-sep}.

Conversely, suppose \((F,\eps)\)-models always exist.  Write \(G=\bigcup_kF_k\) with \(1\in F_k\subseteq F_{k+1}\) finite, and let \(p_k\) be an \((F_k,1/k)\)-model on a set \(Y_k\).  Replacing \(p_k\) by its diagonal action on many disjoint copies of \(Y_k\) preserves \eqref{eq:model-mult} and \eqref{eq:model-sep} and makes \(\abs{Y_k}\to\infty\).  Extend \(p_k\) to \(\sigma_k:G\to\Sym(Y_k)\) by the identity outside \(F_k\cup F_k^2\).  For fixed \(g,h\in G\), eventually \(g,h\in F_k\), so the multiplication defects tend to \(0\).  Taking \(g=h=1\) in \eqref{eq:model-mult} gives \(\distH(p_k(1),\id)<1/k\) by bi-invariance, and then \eqref{eq:model-sep} with \(h=1\) gives \(\distH(p_k(g),\id)>1-2/k\) for \(1\ne g\in F_k\).  Hence \((\sigma_k)\) is a sofic approximation.
\end{proof}

\begin{theorem}[Finite multiplication-table presentation]\label{thm:finite-table}
Let \(G\) be a finitely generated nonsofic group.  Then there is a finite symmetric set
\[
 1\in F=F^{-1}\subseteq G
\]
containing a generating set of \(G\) such that the finitely presented group
\begin{equation}\label{eq:HF-presentation}
 H_F=
 \left\langle
 x_g\ (g\in F\cup F^2)
 \ \middle|\
 x_1=1,\quad x_gx_h=x_{gh}\ (g,h\in F)
 \right\rangle
\end{equation}
is infinite and nonsofic.
\end{theorem}

\begin{proof}
Since \(G\) is nonsofic, Lemma~\ref{lem:models-equivalence} provides a finite set \(F_1\ni1\) and \(\eps_0>0\) admitting no \((F_1,\eps_0)\)-model.  Enlarge \(F_1\) to a finite symmetric set \(F=F^{-1}\) containing \(1\) and a finite generating set of \(G\); since an \((F,\eps_0)\)-model restricts to an \((F_1,\eps_0)\)-model, no \((F,\eps_0)\)-model exists either.  The presentation \eqref{eq:HF-presentation} is finite.  Evaluation of the labels defines a homomorphism
\begin{equation}\label{eq:HF-surjection}
 \pi:H_F\longrightarrow G,
 \qquad
 \pi(x_g)=g.
\end{equation}
Because \(F\) contains a generating set, \(\pi\) is surjective.  Hence \(H_F\) is infinite.

Suppose that \(H_F\) were sofic.  Choose a sufficiently accurate sofic approximation of a finite subset containing all \(x_g\), \(g\in F\cup F^2\), and all products and inverses needed below.  Write \(p(g)\) for the permutation assigned to \(x_g\).  Since
\[
 x_gx_h=x_{gh}
 \qquad(g,h\in F)
\]
holds exactly in \(H_F\), approximate multiplicativity gives, after choosing the approximation accuracy small enough,
\[
 \distH\bigl(p(g)p(h),p(gh)\bigr)<\eps_0
 \qquad(g,h\in F).
\]
If \(g\ne h\) in \(F\), then \(x_g\ne x_h\), since their images under \(\pi\) are distinct; thus \(w=x_g^{-1}x_h\ne1\) in \(H_F\).  Normalize as in Lemma~\ref{lem:finite-normalization}, so that \(\sigma(x_g^{-1})=\sigma(x_g)^{-1}\).  Bi-invariance of \(\distH\) gives
\[
 \distH\bigl(p(g),p(h)\bigr)
 =\distH\bigl(\id,\sigma(x_g)^{-1}\sigma(x_h)\bigr)
 \geq
 \distH\bigl(\id,\sigma(w)\bigr)
 -\distH\bigl(\sigma(x_g^{-1})\sigma(x_h),\sigma(w)\bigr).
\]
Sofic faithfulness makes the first term at least \(1-\eps_0/2\), and approximate multiplicativity at the pair \((x_g^{-1},x_h)\), whose product is \(x_g^{-1}x_h=w\), makes the second term at most \(\eps_0/2\), for a sufficiently accurate approximation.  Hence \(\distH(p(g),p(h))>1-\eps_0\), and \(p\) is a forbidden \((F,\eps_0)\)-model of \(G\), a contradiction.
\end{proof}

\begin{theorem}[Kazhdan cover]\label{thm:kazhdan-cover}
Every property-\((T)\) group that is not sofic is a quotient of an infinite, finitely presented property-\((T)\) group that is not sofic.
\end{theorem}

\begin{proof}
A discrete property-\((T)\) group \(G\) is finitely generated.  Choose \(F\) and \(\eps_0\) as in the proof of Theorem~\ref{thm:finite-table}.  By Shalom's theorem \cite{Shalom00}, every discrete group with property~\((T)\) is a quotient of a finitely presented group with property~\((T)\); fix such a group \(P\) and a surjection \(\rho:P\onto G\).  For \(g\in F\cup F^2\), choose lifts \(\widetilde g\in P\) with \(\rho(\widetilde g)=g\) and \(\widetilde 1=1\), and put
\[
 H=P\Big/\bigl\langle\!\bigl\langle
 \widetilde g\,\widetilde h\,\widetilde{gh}^{\,-1}
 :g,h\in F
 \bigr\rangle\!\bigr\rangle .
\]
Every added relator lies in \(\ker\rho\), so \(\rho\) factors through a surjection \(H\onto G\); hence \(H\) is infinite, and the images \(x_g\in H\) of the lifts are distinct for distinct \(g\in F\cup F^2\).  The group \(H\) is finitely presented, and it has property~\((T)\) as a quotient of \(P\).  The relations \(x_gx_h=x_{gh}\) for \(g,h\in F\), together with \(x_1=1\), hold exactly in \(H\), so the argument of Theorem~\ref{thm:finite-table}, applied verbatim to the elements \(x_g\), shows that a sofic approximation of \(H\) would produce an \((F,\eps_0)\)-model of \(G\).  Hence \(H\) is not sofic.
\end{proof}

\begin{remark}[Effectivity]
Theorems~\ref{thm:finite-table} and~\ref{thm:kazhdan-cover} are existential: nonsoficity supplies the finite obstruction set \(F\) and the threshold \(\eps_0\), but the proof does not compute them.  They produce finite presentations in the logical sense, not a printed list of relators.
\end{remark}

\begin{proof}[Proof of Theorem~\ref{thm:main}]
Corollary~\ref{cor:G-properties} proves the first assertion.  Apply Theorem~\ref{thm:finite-table} to that group to obtain \(H_F\), and Theorem~\ref{thm:kazhdan-cover} for the property-\((T)\) refinement.
\end{proof}

\section{Concluding remarks}\label{sec:remarks}

\begin{remark}[Elementary memberships]
The explicit factorizations in Lemma~\ref{lem:two-by-two}, Proposition~\ref{prop:u-elementary}, Lemma~\ref{lem:z-properties}, and Proposition~\ref{prop:J-in-Gamma} make the argument independent of the equality \(\GL_m(R)=\EL_m(R)\), and hence of \(K_1(R)=0\): no \(K\)-theory is used.  The argument says nothing about hyperlinearity or the Connes embedding problem.
\end{remark}

\begin{remark}[Scope]
Nothing in Sections~\ref{sec:compressors}--\ref{sec:nonsofic} depends on the rank four: the same matrices, written in rank \(m+1\) over the left comb with \(m+1\) leaves, show that \(\EL_{m+1}(L_{\F_2}(1,2))\) is not sofic for every \(m\geq3\).  Over an arbitrary field \(k\), the construction applies whenever the matrices \(u\), \(z\), and \(j(g)\) are known to lie in the relevant elementary groups; this follows, for instance, from the vanishing of \(K_1(L_k(1,2))\) and the GE property of purely infinite simple rings, see \cite{AbramsAranda06,AraGoodearlPardo02,AraBrustengaCortinas09}.  A leaf-set isomorphism \(M_r(L)\cong L\) then propagates nonsoficity to the full unit group \(L^\times\) and to every \(\GL_r(L)\).  We do not pursue this here.
\end{remark}

\begin{remark}
The two compressors have different roles.  Analytically, one compressor \(q_0\) is enough to produce the matched component carrying the expanding \(K\)-action.  Algebraically, the pair is useful because
\[
 vu^{-1}=z
\]
is the block involution that generates all elementary matrices involving the fourth coordinate.
\end{remark}

\begin{remark}
The value \(1/2\) in the median argument is a normalization, not a corner trace.  It records equality with the median size:
\[
 \frac{s}{s+m}=\frac12
 \quad\Longleftrightarrow\quad s=m.
\]
The conservation identity is what converts the one-sided component comparison into two-sided size control.
\end{remark}

\begin{remark}
The passage from \(G\) to \(H_F\) does not use closure of soficity under quotients.  The surjection \(H_F\onto G\) is used only to prove infinitude and to distinguish the finitely many named generators.  Nonsoficity is inherited from the explicit forbidden finite-table test.
\end{remark}

\begin{thebibliography}{99}

\bibitem{AbramsAranda06}
G.~Abrams and G.~Aranda Pino,
\emph{Purely infinite simple Leavitt path algebras},
J. Pure Appl. Algebra \textbf{207} (2006), 553--563.

\bibitem{AraBrustengaCortinas09}
P.~Ara, M.~Brustenga, and G.~Corti\~nas,
\emph{$K$-theory of Leavitt path algebras},
M\"unster J. Math. \textbf{2} (2009), 5--34.

\bibitem{AraGoodearlPardo02}
P.~Ara, K.~R. Goodearl, and E.~Pardo,
\emph{$K_0$ of purely infinite simple regular rings},
$K$-Theory \textbf{26} (2002), 69--100.

\bibitem{BaumslagSolitar62}
G.~Baumslag and D.~Solitar,
\emph{Some two-generator one-relator non-Hopfian groups},
Bull. Amer. Math. Soc. \textbf{68} (1962), 199--201.

\bibitem{CannonFloydParry96}
J.~W. Cannon, W.~J. Floyd, and W.~R. Parry,
\emph{Introductory notes on Richard Thompson's groups},
Enseign. Math. (2) \textbf{42} (1996), 215--256.

\bibitem{BleakQuick17}
C.~Bleak and M.~Quick,
\emph{The infinite simple group $V$ of Richard J. Thompson: presentations by permutations},
Groups Geom. Dyn. \textbf{11} (2017), 1401--1436.

\bibitem{ElekSzabo06}
G.~Elek and E.~Szab\'o,
\emph{On sofic groups},
J. Group Theory \textbf{9} (2006), 161--171.

\bibitem{ErshovJaikin10}
M.~Ershov and A.~Jaikin-Zapirain,
\emph{Property $(T)$ for noncommutative universal lattices},
Invent. Math. \textbf{179} (2010), 303--347.

\bibitem{Gromov99}
M.~Gromov,
\emph{Endomorphisms of symbolic algebraic varieties},
J. Eur. Math. Soc. \textbf{1} (1999), 109--197.

\bibitem{Higman74}
G.~Higman,
\emph{Finitely Presented Infinite Simple Groups},
Notes on Pure Mathematics, vol.~8,
Department of Pure Mathematics, Australian National University, Canberra, 1974.

\bibitem{Kun19}
G.~Kun,
\emph{On sofic approximations of property $(T)$ groups},
arXiv:1606.04471v5, 2019.

\bibitem{KunThom19}
G.~Kun and A.~Thom,
\emph{Inapproximability of actions and Kazhdan's property $(T)$},
arXiv:1901.03963v2, 2019.

\bibitem{Leavitt62}
W.~G. Leavitt,
\emph{The module type of a ring},
Trans. Amer. Math. Soc. \textbf{103} (1962), 113--130.

\bibitem{Nekrashevych04}
V.~Nekrashevych,
\emph{Cuntz--Pimsner algebras of group actions},
J. Operator Theory \textbf{52} (2004), 223--249.

\bibitem{Pardo11}
E.~Pardo,
\emph{The isomorphism problem for Higman--Thompson groups},
J. Algebra \textbf{344} (2011), 172--183.

\bibitem{Pestov08}
V.~Pestov,
\emph{Hyperlinear and sofic groups: a brief guide},
Bull. Symbolic Logic \textbf{14} (2008), 449--480.

\bibitem{Shalom00}
Y.~Shalom,
\emph{Rigidity of commensurators and irreducible lattices},
Invent. Math. \textbf{141} (2000), 1--54.

\bibitem{VershikGordon97}
A.~M. Vershik and E.~I. Gordon,
\emph{Groups that are locally embeddable in the class of finite groups},
Algebra i Analiz \textbf{9} (1997), 71--97; English transl., St.~Petersburg Math. J. \textbf{9} (1998), 49--67.

\bibitem{Weiss00}
B.~Weiss,
\emph{Sofic groups and dynamical systems},
Sankhy\=a Ser. A \textbf{62} (2000), 350--359.

\end{thebibliography}

\end{document}
